\PassOptionsToPackage{dvipsnames}{xcolor}
\documentclass[14pt,a4paper,oneside]{extreport}


\usepackage{cmap}
\usepackage[T2A]{fontenc}
\usepackage[utf8]{inputenc}
\usepackage[english,russian]{babel}
\usepackage{enumitem}
\usepackage{indentfirst}
\usepackage{amsmath,amsfonts,,amssymb,amsthm,mathtools}
\theoremstyle{definition}
\newtheorem{theorem}{Теорема}[section]
\newtheorem{corollary}[theorem]{Следствие}

\DeclareMathOperator{\DD}{d}

\usepackage{siunitx}
\sisetup{output-decimal-marker={,}}

\usepackage[version=4]{mhchem}

\usepackage{hyperref}
\hypersetup{
    colorlinks=true,
    linkcolor=blue!95!black,
    filecolor=black,
    urlcolor=black,
    citecolor=black,
    pdftitle={Курс современной школьной физики},
    pdfpagemode=UseOutlines,
}
\urlstyle{same}

\usepackage{graphicx,xcolor}
\usepackage{subcaption}
\graphicspath{figures}
\setlength\fboxsep{3pt}
\setlength\fboxrule{1pt}
\usepackage{wrapfig}
\usepackage{float}
\usepackage{booktabs}
\usepackage{caption}
\captionsetup{
  font=small
}
\usepackage{anyfontsize}
\captionsetup{justification=centering}


\usepackage{geometry}
\geometry{
  a4paper,
  top=20mm,
  bottom=25mm,
  left=15mm,
  right=15mm
}
\newcommand{\name}[1]{%
  \begin{center}
    \vspace*{0.1cm}
    {\bfseries\fontsize{28}{32}\selectfont
      \parbox{\textwidth}{\centering #1\par}}%
  \end{center}%
}

\newcommand{\DEF}[2]{%
  \refstepcounter{definition}%
  \begin{DefinitionBox}
    \textbf{#1} — #2
  \end{DefinitionBox}
}

\newcommand{\opr}[3]{%
  \refstepcounter{definition}%
  \begin{DefinitionBox}
    \textbf{#1} $\equiv #2$ — #3
  \end{DefinitionBox}
}

\newcommand{\Opr}[4]{%
  \refstepcounter{definition}%
  \begin{DefinitionBox}
    \textbf{#1} $\equiv #2$ — #3; \quad $\left[#2\right] = #4$
  \end{DefinitionBox}
}

\newcommand{\Oprf}[4]{%
  \refstepcounter{definition}%
  \begin{DefinitionBox}
    \textbf{#1} $\equiv #2$ — #3;\\
    #4
  \end{DefinitionBox}
}


\newcommand{\DEFlabel}[3]{%
  \refstepcounter{definition}%
  \begin{DefinitionBox}[title={Определение \thedefinition}]%
    \textbf{#1} — #2%
    \label{#3}%
  \end{DefinitionBox}%
}

\newcommand{\zm}[1]{%
  \begin{RemarkBox}
    #1
  \end{RemarkBox}
}


\newcounter{definition}[chapter]
\renewcommand{\thedefinition}{\thechapter.\arabic{definition}}


\usepackage{tikz}
\usetikzlibrary{arrows.meta,calc,positioning,mindmap,shapes.symbols}

\usepackage{pgfplots}
\pgfplotsset{compat=newest}

\usepackage[most]{tcolorbox}
\tcbuselibrary{breakable, skins}

\newtcolorbox{DefinitionBox}[1][]{
    colback=brown!5!white,
    colframe=orange!80!black,
    fonttitle=\bfseries,
    title={Определение},
    breakable,
    enhanced,
    sharp corners=south,
    boxrule=0.5pt,
    #1
}


\newtcolorbox{RemarkBox}[1][]{
  colback=gray!10,
  colframe=gray,
  fonttitle=\bfseries,
  title={Замечание},
  #1,
  breakable
}

\newtcolorbox{ExampleBox}[1][]{
  colback=green!5!white,
  colframe=green!50!black,
  fonttitle=\bfseries,
  title={Пример},
  breakable,
  #1
}

\usetikzlibrary{decorations.markings,decorations.pathreplacing}

\newcommand{\TikzProvodnikOne}{%
\begin{tikzpicture}[scale=1,>=latex]
  \foreach \y in {-2,-1,0,1,2} {
    \draw[
      postaction={decorate},
      decoration={markings, mark=at position 1 with {\arrow{>}}}
    ] (0,\y) -- (6,\y);
  }
\node at (4.9,2.35) {$\vec{E}_0$};

\filldraw[fill=gray, fill opacity=0.2, draw=black, thick] (3,0) circle (1.5);

\foreach \ang in {60,20,-20,-60} {
  \node[red] at ($(3,0)+(\ang:1.7)$) {$+$};
}
\foreach \ang in {120,160,200,240} {
  \node[blue] at ($(3,0)+(\ang:1.7)$) {$-$};
}

\draw[<-, magenta, semithick] (1.7,0) -- (4.3,0);
\draw[<-, magenta, semithick] (2.1,-1.1) -- (3.8,-1.1);
\draw[<-, magenta, semithick] (2.1,1.1) -- (3.8,1.1);
\node[above] at (3,0) {$\vec{E}_i$};

\node[above] at (3,-1) {$\vec{E}=\vec{0}$};
\end{tikzpicture}%
}

\newcommand{\TikzProvodnikTwo}{%
\begin{tikzpicture}[scale=1,>=latex]
\foreach \ang in {60,20,-20,-60} {
  \node[red] at ($(3,0)+(\ang:1.7)$) {$+$};
}
\foreach \ang in {120,160,200,240} {
  \node[blue] at ($(3,0)+(\ang:1.7)$) {$-$};
}

\begin{scope}[shift={(3,3)}]
  \draw[MidnightBlue, semithick,
    postaction={decorate},
    decoration={markings,
      mark=at position 0.15 with {\arrow{>}},
      mark=at position 0.85 with {\arrow{>}}}
  ]
    (-4,0) .. controls (-2,0) and (-1,-0.3) .. (0,-0.3)
          .. controls ( 1,-0.3) and ( 2,0) .. (4,0);
\end{scope}

\begin{scope}[shift={(3,2.5)}]
  \draw[MidnightBlue, semithick,
    postaction={decorate},
    decoration={markings,
      mark=at position 0.15 with {\arrow{>}},
      mark=at position 0.85 with {\arrow{>}}}
  ]
    (-4,0) .. controls (-2,0) and (-1,-0.6) .. (0,-0.6)
          .. controls ( 1,-0.6) and ( 2,0) .. (4,0);
\end{scope}

\begin{scope}[shift={(3,2)}]
  \draw[MidnightBlue, semithick,
    postaction={decorate},
    decoration={markings,
      mark=at position 0.15 with {\arrow{>}},
      mark=at position 0.85 with {\arrow{>}}}
  ]
    (-4,0) .. controls (-2,0) and (-1,-1.5) .. (0,-1.5)
          .. controls ( 1,-1.5) and ( 2,0) .. (4,0);
\end{scope}

\begin{scope}[shift={(3,-3)}]
  \draw[MidnightBlue, semithick,
    postaction={decorate},
    decoration={markings,
      mark=at position 0.15 with {\arrow{>}},
      mark=at position 0.85 with {\arrow{>}}}
  ]
    (-4,0) .. controls (-2,0) and (-1,0.3) .. (0,0.3)
          .. controls ( 1,0.3) and ( 2,0) .. (4,0);
\end{scope}

\begin{scope}[shift={(3,-2.5)}]
  \draw[MidnightBlue, semithick,
    postaction={decorate},
    decoration={markings,
      mark=at position 0.15 with {\arrow{>}},
      mark=at position 0.85 with {\arrow{>}}}
  ]
    (-4,0) .. controls (-2,0) and (-1,0.6) .. (0,0.6)
          .. controls ( 1,0.6) and ( 2,0) .. (4,0);
\end{scope}

\begin{scope}[shift={(3,-2)}]
  \draw[MidnightBlue, semithick,
    postaction={decorate},
    decoration={markings,
      mark=at position 0.15 with {\arrow{>}},
      mark=at position 0.85 with {\arrow{>}}}
  ]
    (-4,0) .. controls (-2,0) and (-1,1.5) .. (0,1.5)
          .. controls ( 1,1.5) and ( 2,0) .. (4,0);
\end{scope}

\node[above] at (5.7,0) {$\vec{E}$};

\draw[MidnightBlue,semithick,
  postaction={decorate},
  decoration={markings,
    mark=at position 0.15 with {\arrow{>}},
    mark=at position 0.85 with {\arrow{>}}}
] (-1,0) -- (7,0);

\filldraw[fill=white, fill opacity=1, draw=black, thick] (3,0) circle (1.5);
\filldraw[fill=gray, fill opacity=0.2, draw=black, thick] (3,0) circle (1.5);
\node at (3,0) {$\varphi_\text{пр}=\text{const}$};
\end{tikzpicture}%
}

\newcommand{\TikzProvodnikThreeA}{%
\begin{tikzpicture}[scale=1,>=latex]
\foreach \ang in {60,20,-20,-60} {
  \node[red] at ($(3,0)+(\ang:1.7)$) {$+$};
}
\foreach \ang in {120,160,200,240} {
  \node[blue] at ($(3,0)+(\ang:1.7)$) {$-$};
}

\begin{scope}[shift={(3,3)}]
  \draw[MidnightBlue, semithick, postaction={decorate}, decoration={markings, mark=at position 0.15 with {\arrow{>}}, mark=at position 0.85 with {\arrow{>}}}]
    (-3.5,0) .. controls (-2,0) and (-1,-0.3) .. (0,-0.3) .. controls (1,-0.3) and (2,0) .. (3.5,0);
\end{scope}
\begin{scope}[shift={(3,2.5)}]
  \draw[MidnightBlue, semithick, postaction={decorate}, decoration={markings, mark=at position 0.15 with {\arrow{>}}, mark=at position 0.85 with {\arrow{>}}}]
    (-3.5,0) .. controls (-2,0) and (-1,-0.6) .. (0,-0.6) .. controls (1,-0.6) and (2,0) .. (3.5,0);
\end{scope}
\begin{scope}[shift={(3,2)}]
  \draw[MidnightBlue, semithick, postaction={decorate}, decoration={markings, mark=at position 0.15 with {\arrow{>}}, mark=at position 0.85 with {\arrow{>}}}]
    (-3.5,0) .. controls (-2,0) and (-1,-1.5) .. (0,-1.5) .. controls (1,-1.5) and (2,0) .. (3.5,0);
\end{scope}
\begin{scope}[shift={(3,-3)}]
  \draw[MidnightBlue, semithick, postaction={decorate}, decoration={markings, mark=at position 0.15 with {\arrow{>}}, mark=at position 0.85 with {\arrow{>}}}]
    (-3.5,0) .. controls (-2,0) and (-1,0.3) .. (0,0.3) .. controls (1,0.3) and (2,0) .. (3.5,0);
\end{scope}
\begin{scope}[shift={(3,-2.5)}]
  \draw[MidnightBlue, semithick, postaction={decorate}, decoration={markings, mark=at position 0.15 with {\arrow{>}}, mark=at position 0.85 with {\arrow{>}}}]
    (-3.5,0) .. controls (-2,0) and (-1,0.6) .. (0,0.6) .. controls (1,0.6) and (2,0) .. (3.5,0);
\end{scope}
\begin{scope}[shift={(3,-2)}]
  \draw[MidnightBlue, semithick, postaction={decorate}, decoration={markings, mark=at position 0.15 with {\arrow{>}}, mark=at position 0.85 with {\arrow{>}}}]
    (-3.5,0) .. controls (-2,0) and (-1,1.5) .. (0,1.5) .. controls (1,1.5) and (2,0) .. (3.5,0);
\end{scope}

\node[above] at (5.7,0) {$\vec{E}$};
\draw[MidnightBlue,semithick, postaction={decorate}, decoration={markings, mark=at position 0.15 with {\arrow{>}}, mark=at position 0.85 with {\arrow{>}}}] (-0.5,0) -- (6.5,0);
\filldraw[fill=white, fill opacity=1, draw=black, thick] (3,0) circle (1.5);
\filldraw[fill=gray, fill opacity=0.2, draw=black, thick] (3,0) circle (1.5);
\node at (3,0) {$\vec{E}_\text{внутр}=\vec{0}$};
\end{tikzpicture}%
}

\newcommand{\TikzProvodnikThreeB}{%
\begin{tikzpicture}[scale=1,>=latex]
\foreach \ang in {60,20,-20,-60} {
  \node[red] at ($(3,0)+(\ang:1.7)$) {$+$};
}
\foreach \ang in {120,160,200,240} {
  \node[blue] at ($(3,0)+(\ang:1.7)$) {$-$};
}

\begin{scope}[shift={(3,3)}]
  \draw[MidnightBlue, semithick, postaction={decorate}, decoration={markings, mark=at position 0.15 with {\arrow{>}}, mark=at position 0.85 with {\arrow{>}}}]
    (-3.5,0) .. controls (-2,0) and (-1,-0.3) .. (0,-0.3) .. controls (1,-0.3) and (2,0) .. (3.5,0);
\end{scope}
\begin{scope}[shift={(3,2.5)}]
  \draw[MidnightBlue, semithick, postaction={decorate}, decoration={markings, mark=at position 0.15 with {\arrow{>}}, mark=at position 0.85 with {\arrow{>}}}]
    (-3.5,0) .. controls (-2,0) and (-1,-0.6) .. (0,-0.6) .. controls (1,-0.6) and (2,0) .. (3.5,0);
\end{scope}
\begin{scope}[shift={(3,2)}]
  \draw[MidnightBlue, semithick, postaction={decorate}, decoration={markings, mark=at position 0.15 with {\arrow{>}}, mark=at position 0.85 with {\arrow{>}}}]
    (-3.5,0) .. controls (-2,0) and (-1,-1.5) .. (0,-1.5) .. controls (1,-1.5) and (2,0) .. (3.5,0);
\end{scope}
\begin{scope}[shift={(3,-3)}]
  \draw[MidnightBlue, semithick, postaction={decorate}, decoration={markings, mark=at position 0.15 with {\arrow{>}}, mark=at position 0.85 with {\arrow{>}}}]
    (-3.5,0) .. controls (-2,0) and (-1,0.3) .. (0,0.3) .. controls (1,0.3) and (2,0) .. (3.5,0);
\end{scope}
\begin{scope}[shift={(3,-2.5)}]
  \draw[MidnightBlue, semithick, postaction={decorate}, decoration={markings, mark=at position 0.15 with {\arrow{>}}, mark=at position 0.85 with {\arrow{>}}}]
    (-3.5,0) .. controls (-2,0) and (-1,0.6) .. (0,0.6) .. controls (1,0.6) and (2,0) .. (3.5,0);
\end{scope}
\begin{scope}[shift={(3,-2)}]
  \draw[MidnightBlue, semithick, postaction={decorate}, decoration={markings, mark=at position 0.15 with {\arrow{>}}, mark=at position 0.85 with {\arrow{>}}}]
    (-3.5,0) .. controls (-2,0) and (-1,1.5) .. (0,1.5) .. controls (1,1.5) and (2,0) .. (3.5,0);
\end{scope}

\node[above] at (5.7,0) {$\vec{E}$};
\draw[MidnightBlue,semithick, postaction={decorate}, decoration={markings, mark=at position 0.15 with {\arrow{>}}, mark=at position 0.85 with {\arrow{>}}}] (-0.5,0) -- (6.5,0);
\filldraw[fill=white, fill opacity=1, draw=black, thick] (3,0) circle (1.5);
\filldraw[fill=gray, fill opacity=0.2, draw=black, thick] (3,0) circle (1.5);
\filldraw[fill=white, fill opacity=1, draw=black, thick] (3.25,-0.5) circle (0.8);
\node at (3,0.8) {$\vec{E}_\text{внутр}=\vec{0}$};
\node at (3.25,-0.5) {$\vec{E}_\text{пол}=\vec{0}$};
\end{tikzpicture}%
}

\newcommand{\TikzBallA}{%
\begin{tikzpicture}[scale=1,>=latex]
\foreach \ang in {22.5,67.5,...,360} {
  \node[red] at ($(0,0)+(\ang:1.2)$) {$+$};
}
\foreach \ang in {0,45,...,360} {
\draw[MidnightBlue, semithick, postaction={decorate}, decoration={markings, mark=at position 0.65 with {\arrow{latex}}}] (\ang:1) -- (\ang:3.5);
}
\node[above=4pt] at (0,0) {$q>0$};
\fill (0,0) circle (1pt);
\draw (0:0) -- (-90:1);
\node at (-65:0.5) {R};
\node[right] at (0,2.55) {$\vec{E}$};
\filldraw[fill=gray, fill opacity=0.2, draw=black, thick] (0,0) circle (1);
\node[] at (165:3) {$\varepsilon$};
\end{tikzpicture}%
}

\newcommand{\TikzBallB}{%
\begin{tikzpicture}[scale=1,>=latex]
\foreach \ang in {22.5,67.5,...,360} {
  \node[blue] at ($(0,0)+(\ang:1.2)$) {$-$};
}
\foreach \ang in {0,45,...,360} {
\draw[MidnightBlue, semithick, postaction={decorate}, decoration={markings, mark=at position 0.35 with {\arrow{latex}}}] (\ang:3.5) -- (\ang:1);
}
\node[above=4pt] at (0,0) {$q<0$};
\fill (0,0) circle (1pt);
\draw (0:0) -- (-90:1);
\node at (-65:0.5) {R};
\node[right] at (0,2.75) {$\vec{E}$};
\filldraw[fill=gray, fill opacity=0.2, draw=black, thick] (0,0) circle (1);
\node[] at (165:3) {$\varepsilon$};
\end{tikzpicture}%
}

\newcommand{\TikzGraphField}{%
\begingroup
\def\R{2}
\def\A{8}
\begin{tikzpicture}
\begin{axis}[clip=false,
  axis lines=middle,
  xlabel={$r$},
  ylabel={$E$},
  xlabel style={anchor=north},
  ylabel style={anchor=east},
  xmin=0, xmax=6,
  ymin=0, ymax=2.5,
  xticklabels={$R$},
  yticklabels={$k\,\frac{|q|}{\varepsilon R^2}$},
  xtick={\R},
  ytick={{\A/(\R*\R)}},
  samples=200,
]
\addplot[very thick, Green] coordinates {(0,0) (\R,0)};
\addplot[very thick, Green, domain=\R:6] {\A/x^2};
\node[above=3pt, font=\scriptsize] at (axis cs:4,{\A/(4*4)}) {$\propto \frac{1}{r^2}$};
\addplot[densely dashed] coordinates {(\R,0) (\R,{\A/(\R*\R)})};
\addplot[densely dashed] coordinates {(0,{\A/(\R*\R)}) (\R,{\A/(\R*\R)})};
\addplot[only marks, mark=*, mark options={fill=white, draw=black}] coordinates {(\R,{\A/(\R*\R)})};
\addplot[only marks, mark=*, mark options={fill=white, draw=black}] coordinates {(\R,0)};
\node[anchor=north east] at (axis cs:0,0) {$0$};
\node[above, font=\scriptsize] at (axis cs:1,0) {$E_\text{внутр}=0$};
\end{axis}
\end{tikzpicture}
\endgroup%
}

\newcommand{\TikzDielOne}{%
\begin{tikzpicture}[scale=1,>=latex]
\foreach \y in {-2,-1,0,1,2} {
  \draw[
    postaction={decorate},
    decoration={markings, mark=at position 1 with {\arrow{>}}}
  ] (-3,\y) -- (3,\y);
}
\node at (2.9,2.35) {$\vec{E}_0$};
\filldraw[fill=Apricot, fill opacity=0.2, draw=black,thick] (-2,-1.5) rectangle (2,1.5);
\foreach \y in {-0.5,0.5} {
  \draw[semithick, RoyalPurple,
    postaction={decorate},
    decoration={markings, mark=at position 0.02 with {\arrow{<}}}
  ] (-1.7,\y) -- (1.7,\y);
}
\node[above=-2pt] at (0,0.5) {$\vec{E}_i$};
\foreach \y in {-1.5,-0.5,...,1.5} {
  \node[red] at (2.2,\y) {$+$};
}
\foreach \y in {-1.5,-0.5,...,1.5} {
  \node[blue] at (-2.2,\y) {$-$};
}
\end{tikzpicture}%
}

\newcommand{\TikzDielTwo}{%
\begin{tikzpicture}[scale=1,>=latex]
\filldraw[fill=Apricot, fill opacity=0.2, draw=black,thick] (-2,-1.5) rectangle (2,1.5);
\foreach \y in {-1.5,-0.5,...,1.5} {
  \node[red] at (2.2,\y) {$+$};
}
\foreach \y in {-1.5,-0.5,...,1.5} {
  \node[blue] at (-2.2,\y) {$-$};
}
\draw[->, semithick] (0,0.2) -- (1.8,0.2) node[midway, above] {$\vec{E}_0$};
\draw[->, semithick, RoyalPurple] (0,0.2) -- (-1,0.2) node[midway, above] {$\vec{E}_i$};
\draw[->, thick, MidnightBlue] (0,0) -- (0.8,0) node[midway, below] {$\vec{E}$};
\fill (0,0.2) circle (0.04);
\end{tikzpicture}%
}

\newcommand{\TikzPolDiel}{%
\begin{tikzpicture}[scale=1,>=latex]
\draw[gray,semithick] (-1,0) -- (1,0);
\fill[red] (-1,0) circle (4pt);
\fill[blue] (1,0) circle (4pt);
\node[below=4pt] at (-1,0) {$+\,q$};
\node[below=4pt] at (1,0) {$-\,q$};
\end{tikzpicture}%
}

\newcommand{\TikzPolDielNoField}{%
\begingroup
\begin{tikzpicture}[scale=1,>=latex]
\pgfmathsetseed{2026}
\def\nx{14}
\def\ny{6}
\def\sx{1}
\def\sy{1}
\foreach \j in {0,...,\numexpr\ny-1\relax}{
  \foreach \i in {0,...,\numexpr\nx-1\relax}{
    \pgfmathsetmacro{\ang}{360*rnd}
    \begin{scope}[shift={(\i*\sx,\j*\sy)}, rotate=\ang]
      \draw[gray] (-0.25,0) -- (0.25,0);
      \fill[red] (-0.25,0) circle (2.5pt);
      \fill[blue] (0.25,0) circle (2.5pt);
    \end{scope}
  }
}
\end{tikzpicture}
\endgroup%
}

\newcommand{\TikzPolDielField}{%
\begingroup
\begin{tikzpicture}[scale=1,>=latex]
\pgfmathsetseed{2026}
\def\nx{14}
\def\ny{6}
\def\sx{1}
\def\sy{1}
\foreach \j in {0,...,\numexpr\ny-1\relax}{
  \foreach \i in {0,...,\numexpr\nx-1\relax}{
    \pgfmathsetmacro{\ang}{-180 + 30*(2*rnd-1)}
    \begin{scope}[shift={(\i*\sx,\j*\sy)}, rotate=\ang]
      \draw[gray] (-0.25,0) -- (0.25,0);
      \fill[red] (-0.25,0) circle (2.5pt);
      \fill[blue] (0.25,0) circle (2.5pt);
    \end{scope}
  }
}
\draw[semithick, postaction={decorate}, decoration={markings, mark=at position 1 with {\arrow{>}}}] (-2,-0.5) -- (\nx+1,-0.5);
\draw[semithick, postaction={decorate}, decoration={markings, mark=at position 1 with {\arrow{>}}}] (-2,2.5) -- (\nx+1,2.5);
\draw[semithick, postaction={decorate}, decoration={markings, mark=at position 1 with {\arrow{>}}}] (-2,5.5) -- (\nx+1,5.5);
\node[above, font=\Large] at (\nx+1,2.5) {$\vec{E}_0$};
\foreach \y in {0,...,\numexpr\ny-1\relax}{
  \node[blue,font=\Large] at (-1,\y) {$-$};
}
\foreach \y in {0,...,\numexpr\ny-1\relax}{
  \node[red,font=\Large] at (\nx,\y) {$+$};
}
\end{tikzpicture}
\endgroup%
}

\newcommand{\TikzAtomNoField}{%
\begin{tikzpicture}[scale=1,>=latex]
\filldraw[fill=red, draw=black] (0,0) circle (0.5);
\node at (0,0) {$+$};
\draw[dashed, gray] (0,0) circle (2.5);
\filldraw[fill=blue,draw=black] (157:2.5) circle (0.2);
\node[white] at (157:2.5) {$-$};
\end{tikzpicture}%
}

\newcommand{\TikzAtomField}{%
\begin{tikzpicture}[scale=1,>=latex]
\draw[postaction={decorate}, decoration={markings, mark=at position 1 with {\arrow{>}}}] (-4,0) -- (4,0);
\draw[postaction={decorate}, decoration={markings, mark=at position 1 with {\arrow{>}}}] (-4,2.5) -- (4,2.5);
\draw[postaction={decorate}, decoration={markings, mark=at position 1 with {\arrow{>}}}] (-4,-2.5) -- (4,-2.5);
\node[above] at (4,2.5) {$\vec{E}_0$};
\filldraw[fill=red, draw=black] (1.5,0) circle (0.5);
\node at (1.5,0) {$+$};
\def\a{3}
\def\b{2}
\pgfmathsetmacro{\c}{sqrt(\a*\a-\b*\b)}
\coordinate (F) at (0,0);
\begin{scope}[shift={(F)}, xshift={-\c}]
\draw[dashed,gray] (0,0) ellipse ({\a} and {\b});
\end{scope}
\filldraw[fill=blue,draw=black] (157:2.8) circle (0.2);
\node[white] at (157:2.8) {$-$};
\end{tikzpicture}%
}

\begin{document}

\pagestyle{empty}
\pagenumbering{gobble}


\name{\S 5. Электростатическое поле в веществе}


Любое вещество состоит из протонов и электронов, взаимодействие которых приводит к образованию устойчивых систем — атомов. Атомы, в свою очередь, также взаимодействуют между собой: при этом взаимодействуют протоны и электроны, принадлежащие различным атомам.

Характер этого взаимодействия может существенно различаться для разных веществ. В результате все вещества можно принципиально разделить на две группы: вещества, в которых имеются свободные носители электрического заряда (\textit{проводники}), и вещества, в которых свободные заряды отсутствуют (\textit{диэлектрики}).

\DEF{Проводник}{вещество, в котором имеются свободные носители электрического заряда, способные перемещаться по всему объёму вещества (\textit{металлы и растворы электролитов (кислоты, щёлочи, основания)}). }

\DEF{Диэлектрик}{вещество, в котором свободные носители электрического заряда отсутствуют, а заряды могут лишь незначительно смещаться относительно своих равновесных положений (\textit{стекло, эбонит, слюда, пластмассы, резина, сухой воздух}).}

До сих пор мы рассматривали электростатическое поле само по себе, не анализируя его воздействие на вещество. Теперь перейдём к рассмотрению влияния электростатического поля на различные среды. Характер этого влияния существенно различается для проводников и диэлектриков.

Помимо этого, ранее было подробно изучено электростатическое поле точечного заряда. Теперь обобщим полученные представления и рассмотрим электростатические поля, создаваемые другими типами источников.

\subsection*{Проводники}

Свободные заряды в проводнике под действием внешнего электростатического поля начинают перераспределяться, в результате чего возникает \textit{явление электростатической индукции}.

\DEF{Электростатическая индукция}{явление перераспределение свободных зарядов в проводнике под действием внешнего электростатического поля.}

\subsubsection*{Поле внутри проводника}

Любая система при постоянных внешних условиях самопроизвольно приходит в состояние термодинамического равновесия, при котором её макропараметры перестают изменяться. Проводник во внешнем электростатическом поле является термодинамической системой, а значит по истечении некоторого (весьма малого) времени \textit{свободные заряды должны перестать перемещаться по проводнику}. 

В чём состоит условие такого равновесия проводника? Чтобы ответить на этот вопрос, обратимся к рисунку~\ref{fig:provodnik-1}, на котором изображён проводящий шар, помещённый во внешнее электростатическое поле $\vec E_0$.

Под действием внешнего поля свободные отрицательные заряды (электроны в металлах) перемещаются против направления поля $\vec E_0$, в то время как в противоположной области остаются положительно заряженные ионы кристаллической решётки. Такое перераспределение зарядов приводит к возникновению индуцированного электростатического поля $\vec E_i$, направленного противоположно внешнему полю $\vec E_0$.
\vspace{-1\baselineskip}
\begin{figure}[H]
  \centering
  \resizebox{0.42\textwidth}{!}{\TikzProvodnikOne}
  \caption{Индуцированное поле $\vec{E}_i$ внутри проводника во внешнем поле $\vec{E}_0$.}
  \label{fig:provodnik-1}
\end{figure}

Перераспределение зарядов продолжается до тех пор, пока электрическая сила $\vec{F}_\text{эл}$,
действующая на свободный заряд $\varDelta q$ стороны результирующего поля $\vec{E}_\text{внутр}$ в любой точке внутри проводника,
не станет равной нулю:
\begin{equation}
\vec F_{\text{эл}} = \varDelta q\cdot\vec E = \vec 0,
\qquad \varDelta q \neq 0
\;\Rightarrow\;
\vec E = \vec 0.
\end{equation}

Напряженность результирующего электростатического поля $\vec{E}_\text{внутр}$ согласно принципу суперпозиции складывается из внешнего $\vec{E}_0$ и индуцированного (наведенного) $\vec{E}_i$ полей:

\begin{equation}
  \vec{E}_\text{внутр}=\vec{E_0}+\vec{E}_i=\vec{0} \;\Rightarrow\;
\vec{E}_i = - \vec E_0.
\end{equation}

Таким образом, в установившемся состоянии внутри проводника индуцированное поле $\vec E_i$ компенсирует внешнее поле $\vec E_0$. 

Приведённые рассуждения остаются справедливыми для произвольного внешнего электростатического поля $\vec E_0$, в том числе и в случае его отсутствия $\vec{E}_0=\vec{0}$. Кроме того, форма и размеры проводника, а также наличие у него собственного заряда не изменяют сделанного вывода: \textit{в установившемся состоянии напряжённость электростатического поля внутри проводника равна нулю}.

\DEF{I свойство проводника}{
в состоянии электростатического равновесия напряжённость
электростатического поля в любой точке внутри проводника равна нулю:
\begin{equation*}
  \vec E_{\text{внутр}} = \vec 0.
\end{equation*}
}

\subsubsection*{Потенциал проводника}

Из I свойства проводника следует, что \textit{потенциал всех точек проводника постоянен}.
В самом деле:
\begin{align}
  U_{\text{внутр}} &= E_{\text{внутр}} \cdot d = 0, \\
  U_{\text{внутр}} &= \varphi_{1,\text{внутр}} - \varphi_{2,\text{внутр}} = 0, \\
                   &\Rightarrow \varphi_{1,\text{внутр}} = \varphi_{2,\text{внутр}}.
\end{align}
Следовательно, электростатический потенциал внутри проводника одинаков во всех его точках и равен потенциалу проводника $\varphi_{\text{пр}}$.

\DEF{II свойство проводника}{
в состоянии электростатического равновесия электростатический потенциал
одинаков во всех точках проводника и равен потенциалу проводника (эквипотенциальный объем):
\begin{equation*}
  \varphi_{\text{внутр}} = \varphi_{\text{пр}} = \text{const}.
\end{equation*}
}

\subsubsection*{Поле снаружи проводника}

Что происходит в области вне проводящего шара? Результирующее электрическое поле $\vec E$ образуется суперпозицией внешнего поля $\vec E_0$ и индуцированного поля $\vec E_i$. Картина линий напряженности этого поля удовлетворяет следующим соображениям:
\begin{enumerate}
\item Искажение внешнего поля тем больше, чем ближе точка наблюдения расположена к поверхности шара. На достаточно больших расстояниях от проводника вклад индуцированного поля становится пренебрежимо малым, и внешнее поле практически не изменяется.
\item Линии напряжённости электрического поля пересекают поверхность проводника под прямым углом, поскольку поверхность проводника является эквипотенциальной.
\end{enumerate}
В результате наложения внешнего $\vec{E}_0$ и индуцированного $\vec{E}_i$ полей картина линий напряжённости результирующего электростатического поля $\vec E$ имеет вид, показанный на рисунке~\ref{fig:provodnik-2}.

\begin{figure}[H]
  \centering
  \resizebox{0.52\textwidth}{!}{\TikzProvodnikTwo}
  \caption{Результирующее искаженное электростатическое поле $\vec{E}$. }
  \label{fig:provodnik-2}
\end{figure}

\DEF{III cвойство проводника}{линии напряжённости результирующего электростатического поля
перпендикулярны поверхности проводника в каждой её точке.}

\subsubsection*{Заряд проводника}

I свойство проводника утверждает, что в состоянии электростатического равновесия напряжённость электрического поля внутри проводника равна нулю. Если бы внутри проводника существовал избыточный заряд, он создавал бы электрическое поле, что противоречит этому свойству. Следовательно, \textit{избыточный заряд в объёме проводника отсутствует и весь заряд распределяется по его поверхности.}

\DEF{IV свойство проводника}{
в состоянии электростатического равновесия весь избыточный (нескомпенсированный) электрический
заряд проводника находится на его поверхности.
}

Что произойдёт с электростатическим полем $\vec{E}$ и распределением заряда $q$ проводника, если внутри него создать полость, изъяв часть вещества? Согласно III свойству проводника весь избыточный электрический заряд $q$, если он имеется, расположен на поверхности проводника. Поскольку полость образована внутри проводника и не содержит зарядов, \textit{заряд остаётся распределённым по внешней поверхности проводника прежним образом.}

Следовательно, электростатическое поле внутри проводника $\vec{E}_\text{внутр}$, включая область полости $\vec{E}_\text{пол}$, по-прежнему равно нулю, а поле вне проводника $\vec{E}$ остаётся таким же, как и до образования полости. 

На рисунке~\ref{fig:provodnik-3} показана картина линий напряжённости результирующего электростатического поля $\vec E$ незаряженного проводника, помещённого во внешнее однородное поле $\vec E_0$. На рисунке~\ref{fig:provodnik-3a} изображён проводник без полости, а на рисунке~\ref{fig:provodnik-3b} — проводник с полостью. Как видно, наличие полости внутри проводника не изменяет электростатическое поле ни внутри проводника, ни вне его.

\begin{figure}[H]
  \centering

  \begin{subfigure}{0.46\textwidth}
    \centering
    \resizebox{\linewidth}{!}{\TikzProvodnikThreeA}
    \caption{Проводник без полости.}
      \label{fig:provodnik-3a}
  \end{subfigure}\hspace{1cm}
  \begin{subfigure}{0.47\textwidth}
    \centering
    \resizebox{\linewidth}{!}{\TikzProvodnikThreeB}
    \caption{Проводник c полостью.}
      \label{fig:provodnik-3b}
  \end{subfigure}
  \caption{Незаряженный проводник во внешнем поле.}
  \label{fig:provodnik-3}
\end{figure}

Явление независимости электростатического поля внутри полости проводника $\vec E_{\text{пол}}$ от внешнего поля $\vec E_0$ называется \textit{экранированием}. Это явление лежит в основе электростатической защиты приборов от воздействия внешнего электрического поля (например, клетки Фарадея).

\subsubsection*{Поле уединённого заряженного проводящего шара (сферы)}

Рассмотрим электростатическое поле $\vec E$, создаваемое неподвижным
проводящим шаром радиуса $R$ с зарядом $q$.

Согласно I свойству проводника напряжённость электростатического поля
в любой точке внутри шара равна нулю:
\begin{equation}
\vec E_{\text{внутр}} = \vec 0.
\end{equation}

Согласно II свойству проводника потенциал всех точек шара одинаков,
$\varphi_{\text{пр}} = \text{const}$. 

Согласно III свойству проводника силовые линии
электростатического поля $\vec E$ снаружи шара в каждой точке
перпендикулярны его поверхности.

Согласно IV свойству проводника весь заряд $q$ расположен на поверхности
шара. Поскольку шар обладает центральной симметрией, заряд $q$
распределяется равномерно по всей его поверхности.

Приведённые рассуждения справедливы и для заряженной \textit{проводящей сферы}: наличие полости не изменяет распределение избыточного заряда и создаваемое им электростатическое поле.

На рисунке~\ref{fig:ball} показана картина линий напряженности электростатического поля уединенного заряженного проводящего шара (сферы). Случай положительного заряда $q>0$ представлен на рисунке~\ref{fig:ball-a}, а случай отрицательного заряда $q<0$ — на рисунке~\ref{fig:ball-b}.
\begin{figure}[H]
  \centering

  \begin{subfigure}{0.48\textwidth}
    \centering
    \resizebox{\linewidth}{!}{\TikzBallA}
    \caption{$q>0$.}
      \label{fig:ball-a}
  \end{subfigure}\hspace{0.5cm}
  \begin{subfigure}{0.48\textwidth}
    \centering
    \resizebox{\linewidth}{!}{\TikzBallB}
    \caption{$q<0$.}
      \label{fig:ball-b}
  \end{subfigure}
  \caption{Электростатическое поле уединённого заряженного \\ проводящего шара (сферы).}
  \label{fig:ball}
\end{figure}

Как видно из рисунка \ref{fig:ball} электростатическое поле $\vec{E}$ уединенного заряженного проводящего шара (сферы) внутри равно нулю $\vec{E}_\text{внутр}=\vec{0}$, а вне шара совпадает с электростатическим полем точечного заряда $q$, расположенного в центре шара. Поэтому зависимость модуля напряжённости электрического поля $E$ от расстояния $r$ до центра шара (сферы) радиуса $R$ имеет вид:
\begin{center}
\begin{minipage}{0.43\textwidth}
\begin{align*}
\boxed{
E(r)=
\begin{cases}
0, & \text{при } r<R,\\[4pt]
\dfrac{k\,|q|}{\varepsilon\,r^2}, & \text{при } r > R.
\end{cases}}
\end{align*}
\end{minipage}
\hfill
\begin{minipage}{0.56\textwidth}
\centering
\resizebox{\linewidth}{!}{\TikzGraphField}
\end{minipage}
\end{center}

\subsection*{Диэлектрики}

Диэлектрики, в отличие от проводников, не содержат свободных носителей электрического заряда. Вследствие этого их поведение в электрическом поле качественно отличается от поведения проводников. Так как все ранее полученные свойства проводника опирались на возможность свободного перераспределения зарядов, они не применимы к диэлектрикам.

При помещении диэлектрика во внешнее электрическое поле $\vec E_0$ на его торцах возникает нескомпенсированный \textit{связанный заряд}, создающий электростатическое поле $\vec E_i$, которое направлено противоположно внешнему полю $\vec{E}_0$ (рис.~\ref{fig:diel-1}).
\vspace{-1\baselineskip}
\begin{figure}[H]
  \centering
  \resizebox{0.45\textwidth}{!}{\TikzDielOne}
  \caption{Индуцированное поле $\vec{E}_i$ внутри диэлектрика во внешнем поле $\vec{E}_0$.}
  \label{fig:diel-1}
\end{figure}

\subsubsection*{Диэлектрическая проницаемость}

Согласно принципу суперпозиции результирующая напряжённость $\vec E$ электростатического поля внутри диэлектрика равна сумме напряжённостей внешнего поля $\vec E_0$ и поля $\vec E_i$, создаваемого индуцированными связанными зарядами:
\begin{equation}
\vec E = \vec E_0 + \vec E_i.
\end{equation}

Поскольку в диэлектрике отсутствуют свободные заряды, поле $\vec E_i$, создаваемое связанными зарядами, не может полностью компенсировать внешнее поле $\vec E_0$. В результате напряжённость результирующего поля $\vec E$ внутри диэлектрика сонаправлена с $\vec E_0$, как показано на рисунке \ref{fig:diel-2}. 
\vspace{-1\baselineskip}
\begin{figure}[H]
  \centering
  \resizebox{0.37\textwidth}{!}{\TikzDielTwo}
  \caption{Принцип суперпозиции: $\vec E=\vec E_0+\vec E_i$, $E_i<E_0$. }
  \label{fig:diel-2}
\end{figure}
\noindent Поэтому модуль напряженности $E$ результирующего поля:
\begin{equation}
  E=E_0-E_i. 
  \label{eq:8}
\end{equation}

В случае \textit{линейных диэлектриков} модуль напряжённости индуцированного поля $E_i$
прямо пропорционален модулю напряжённости результирующего поля $E$:
\begin{equation}
  E_i \propto E \;\Rightarrow\; E_i = \chi E,
  \label{eq:9}
\end{equation}
где $\chi$ — диэлектрическая восприимчивость (табличная безразмерная величина).

Выразим модуль напряжённости $E$ результирующего поля внутри линейного диэлектрика
через модуль напряжённости $E_0$ внешнего поля:
\begin{equation}
 E = E_0 - E_i
\;\xrightarrow{\eqref{eq:9}}\;
E = E_0 - \chi E
\;\xrightarrow\;
E = \frac{E_0}{1+\chi}.
  \label{eq:10}
\end{equation}

Для удобства в уравнении~\eqref{eq:10} вместо множителя $(1+\chi)$ вводят
\textit{диэлектрическую проницаемость} $\varepsilon$
\[
\varepsilon = 1+\chi,
\]
которая показывает, во сколько раз модуль напряжённости электрического поля $E$
в данном веществе меньше модуля напряжённости $E_0$ этого поля в вакууме.

\opr{Диэлектрическая проницаемость вещества}{\varepsilon}{табличная скалярная физическая величина, показывающая во сколько раз данное вещество ослабляет электрическое поле по сравнению с вакуумом: 
\begin{equation*}
\varepsilon = \frac{E_0}{E}, \qquad [\varepsilon]=\si{1}.
\end{equation*}
}

В таблице \ref{tab:epsilon} приведены значения диэлектрической проницаемости некоторых веществ. 

\begin{table}[H]
\centering
\caption{Диэлектрическая проницаемость некоторых веществ}
\label{tab:epsilon}
\renewcommand{\arraystretch}{1.2}
\begin{tabular}{l c @{\hspace{1.5cm}} l c}
\toprule
\textbf{Вещество} & $\boldsymbol{\varepsilon}$ &
\textbf{Вещество} & $\boldsymbol{\varepsilon}$ \\
\midrule
Вакуум                 & 1         & Слюда     & 6 \\
Воздух                 & 1{,}00058 & Стекло    & 7 \\
Керосин                & 2         & Текстолит & 7 \\
Парафин                & 2         & Фарфор    & 5 \\
Масло                  & 2{,}2     & Эбонит    & 3 \\
Парафинированная бумага & 27        & Вода      & 81 \\
\bottomrule
\end{tabular}
\end{table}

Таким образом, вследствие возникновения связанного заряда на поверхности диэлектрика
напряжённость электрического поля $E_0$ уменьшается в $\varepsilon$ раз.
Однако возникает вопрос: откуда появляется этот связанный заряд?
Он обусловлен явлением \textit{поляризации диэлектрика}.

\DEF{Поляризация диэлектрика}{явление смещения связанных электрических зарядов в диэлектрике
под действием внешнего электрического поля.
}

Рассмотрим наиболее распространенные механизмы поляризации:
\textit{дипольную} и \textit{электронную}.

\subsubsection*{Дипольная поляризация}

В молекуле диэлектрика — поваренной соли \ce{NaCl} — внешний электрон атома натрия передаётся атому хлора, которому недостаёт одного электрона для завершения внешней электронной оболочки. В результате атом натрия превращается в положительный ион \ce{Na+}, а атом хлора — в отрицательный ион \ce{Cl-}. Таким образом, молекулу \ce{NaCl} с точки зрения электростатики можно рассматривать как систему, образованную двумя противоположными электрическими зарядами, находящимися на некотором расстоянии друг от друга. Такие системы называют \textit{диполями}, а соответствующие молекулы — \textit{полярными}.

\DEF{Электрический диполь}{система, состоящая из двух противоположных электрических зарядов, находящихся на некотором расстоянии друг от друга.
}

\begin{wrapfigure}{r}{0.38\textwidth}
  \vspace{-20pt}
  \centering
  \resizebox{0.65\linewidth}{!}{\TikzPolDiel}
  \caption{Электрический диполь}
  \label{fig:pol-diel}
\end{wrapfigure}

На рисунке \ref{fig:pol-diel} изображён электрический диполь. К полярным диэлектрикам относятся такие вещества, как вода \ce{H2O}, аммиак \ce{NH3}, хлороводород \ce{HCl} и этанол \ce{C2H5OH}.

В отсутствие внешнего электрического поля ($\vec E_0=\vec 0$) молекулы полярного диэлектрика вследствие теплового движения ориентированы случайным образом, как показано на рисунке \ref{fig:pol-diel-nofield}. В результате суммарный связанный заряд отсутствует, поляризации диэлектрика нет. 

\begin{figure}[H]
  \centering
  \resizebox{0.72\textwidth}{!}{\TikzPolDielNoField}
  \caption{Случайная ориентация полярных молекул в отсутствии внешнего поля ($\vec E_0=\vec{0}$).}
  \label{fig:pol-diel-nofield}
\end{figure}

Во внешнем электрическом поле $\vec E_0$ полярные молекулы диэлектрика приобретают преимущественную ориентацию вдоль направления поля. При этом положительный заряд $+q$ ориентируется по направлению силовых линий, а отрицательный заряд $-q$ — в противоположную сторону, как показано на рисунке \ref{fig:pol-diel-field}. 

\begin{figure}[H]
  \centering
  \resizebox{0.9\textwidth}{!}{\TikzPolDielField}
  \caption{Преимущественная ориентация полярных молекул во внешнем электрическом поле ($\vec E_0 \neq \vec 0$).}
  \label{fig:pol-diel-field}
\end{figure}

Как видно на рисунке \ref{fig:pol-diel-field}, при этом на торцах диэлектрика возникает \textit{связанный нескомпенсированный заряд} (слева — отрицательный, справа — положительный), тогда как в объёме диэлектрика заряды соседних полярных молекул взаимно компенсируются. В результате возникает \textit{дипольная поляризация}. 

\DEF{Дипольная (ориентационная) поляризация}{поляризация, обусловленная преимущественной ориентацией полярных молекул во внешнем электрическом поле.}

\subsubsection*{Электронная поляризация}

Вещество любого диэлектрика состоит из атомов, состоящее согласно модели Резерфорда из положительно заряженного ядра и вокруг вращающихся отрицательно заряженных электронов. 

В отсутствие внешнего электрического поля ($\vec E_0=\vec 0$) (см. рис.~\ref{fig:atom-nofield}) электрон движется по круговой траектории вокруг ядра под действием силы кулоновского притяжения. Его среднее по времени положение совпадает с положением ядра, поэтому поляризация отсутствует.

Под действием внешнего электрического поля $\vec E_0 \neq \vec 0$ круговая орбита электрона искажается и становится эллиптической (см. рис.~\ref{fig:atom-field}). В результате электрон проводит больше времени по одну сторону от ядра, чем по другую, и его среднее (по времени) положение уже не совпадает с положением ядра. Таким образом, в атоме возникает электрический диполь.
\begin{figure}[H]
  \centering

  \begin{subfigure}{0.34\textwidth}
    \centering
    \resizebox{\linewidth}{!}{\TikzAtomNoField}
    \caption{В отсутствие внешнего электрического поля ($\vec E_0 = \vec 0$).}
      \label{fig:atom-nofield}
  \end{subfigure}\hspace{1.5cm}
  \begin{subfigure}{0.55\textwidth}
    \centering
    \resizebox{\linewidth}{!}{\TikzAtomField}
    \caption{Во внешнем электрическом поле ($\vec E_0\neq \vec 0$).}
      \label{fig:atom-field}
  \end{subfigure}
  \caption{Поляризация атома во внешнем электрическом поле (модель Резерфорда).}
  \label{fig:atom}
\end{figure}

\DEF{Электронная поляризация}{поляризация, обусловленная смещением электронных оболочек атомов во внешнем электрическом поле.
}

Разумеется, полярные диэлектрики также подвержены электронной поляризации, однако для них основной вклад в поляризацию вносит дипольная (ориентационная) поляризация. Именно этим объясняется их значительно б\'{о}льшая диэлектрическая проницаемость $\varepsilon$: например, для воды $\varepsilon = 81$ (см. табл.~\ref{tab:epsilon}).


\end{document}
